% s03-1-3-characters.tex —— 3.1.3 小节：B 的主角（印刷页 70–74）

\subsection{B 的主角}
\label{sec:b-characters}

下面回顾 \cite[1.5 节]{ABC+24a} 中把球面 Hecke 范畴 $\Sph_G$ 定义为幺半因子化范畴的
构造。考虑因子化范畴 $D\text{-}\mathrm{mod}_{1/2}(\Gr_G)$，取 $L^+G$-不变量
\[
D\text{-}\mathrm{mod}_{1/2}(\Gr_G)^{L^+G},
\]
它也可写作
\[
D\text{-}\mathrm{mod}_{1/2}(L^+G\backslash LG/L^+G).
\]
这个因子化范畴带有卷积给出的幺半结构。

接着作重整化：考虑 $D\text{-}\mathrm{mod}_{1/2}(L^+G\backslash LG/L^+G)$ 中一切沿
\[
LG/L^+G\to L^+G\backslash LG/L^+G
\]
的 $!$-拉回为紧的对象（等价地，沿 $L^+G\backslash LG\to L^+G\backslash LG/L^+G$ 的
$!$-拉回为紧的对象）。于是

\begin{definition}
\label{def:sph}
球面 Hecke 范畴 $\Sph_G\in \FactAlgCat$ 是上述对象构成的范畴的 Ind-完备化。
它从未重整化的版本继承幺半结构。
\end{definition}
\srcpage{70}

\begin{remark}
\label{rem:sph-twist}
若要让 $\Sph_G$ 作用在 $\Whit^{!}(G)$ 上，则应当改用 $\Sph_{G,\rho(\omega_X)}$：
把 $\rho(\omega_X)$ 经 $T\to G$ 视为主 $G$-丛，并用 \cite[1.2 节]{ABC+24a} 的几何扭转
构造作出
\[
D\text{-}\mathrm{mod}_{1/2}(\Gr_{G,\rho(\omega_X)})^{L^+G^{\rho(\omega_X)},\mathrm{ren}}.
\]
然而由 \cite[1.5.8 节]{ABC+24a}，有典范等同
\[
\Sph_G\simeq \Sph_{G,\rho(\omega_X)},
\]
这来自几何扭转构造的一般事实
\[
Y/L^+H\simeq Y^{P_H}/L^+H^{P_H}.
\]
\end{remark}

下面把仿射 Hecke 范畴 $\Aff_G$ 定义为幺半因子化 $\Sph_G$-模范畴。更精确地说，
对 $(\Sph_G,\Aff_G)$ 将具有 $\FactModCat_{x_0}$ 中幺半对象的结构。

先给出 Iwahori 子群 $I\subset L^+G$ 的因子化模空间版本。回顾 $L^+G_{\Ran}$ 具有因子化群
结构：对试验仿射概形 $S$ 及 $S$-点 $x\colon S\to \Ran$，$x$ 到 $L^+G_{\Ran}$ 的提升
由映射
\[
\widehat{D}_x\to G
\]
（或等价地，因 $G$ 仿射，由映射 $D_x\to G$）给出。
\srcpage{71}

于是 Iwahori $I^{\Ran}_{x_0}$ 是如下定义的因子化 $L^+G^{\Ran}$-模群：对试验仿射概形
$S$ 及 $S$-点 $x\colon S\to \Ran_{x_0}$，$x$ 到 $I^{\Ran}_{x_0}$ 的提升由映射
\[
D_x\to G
\]
给出，并要求 $S\times\{x_0\}\subset D_x$ 落在 $B\subset G$ 中。

接着取因子化 $D\text{-}\mathrm{mod}_{1/2}(\Gr_G)$-模范畴 $D\text{-}\mathrm{mod}_{1/2}(\Fl_G)$，
取 $I$-不变量
\[
D\text{-}\mathrm{mod}_{1/2}(\Fl_G)^I=D\text{-}\mathrm{mod}_{1/2}(I\backslash LG/I),
\]
它是带卷积幺半结构的因子化 $D\text{-}\mathrm{mod}_{1/2}(\Gr_G)^{L^+G}$-模范畴。
按与前面相同的方式重整化，即取其中沿
\[
LG/I\to I\backslash LG/I
\]
的 $!$-拉回为紧的对象（等价地，沿 $I\backslash LG\to I\backslash LG/I$ 的 $!$-拉回为紧
的对象）构成的子范畴再作 Ind-完备化，即得

\begin{definition}
\label{def:aff}
$\FactModCat_{x_0}$ 中的对 $(\Sph_G,\Aff_G)$ 由因子化范畴 $\Sph_G$ 与刚构造的
因子化 $\Sph_G$-模范畴 $\Aff_G$ 构成。对 $(\Sph_G,\Aff_G)$ 从未重整化的版本
（两个分量都如此）继承幺半结构。
\end{definition}

\begin{remark}
\label{rem:aff-twist}
与球面情形一样，为使 $\Aff_G$ 作用在 $\Whit^{!}(\Fl_G)$ 上（更精确地说，为使幺半对
$(\Sph_G,\Aff_G)$ 作用在对 $(\Whit^{!}(G),\Whit^{!}(\Fl_G))$ 上），本应使用经
$\rho(\omega_X)$ 的几何扭转构造。然而与球面情形一样，扭转版本与非扭转版本之间有
典范等同。
\end{remark}
\srcpage{72}

再转到谱侧。谱球面 Hecke 范畴 $\Sph^{\spec}_{\check G}$ 在 \cite[附录 E.2]{ABC+24a}
中定义（附录 E 的其余部分有详细讨论），方式如下。

考虑\emph{谱球面 Hecke 叠}
\[
\Hck^{\spec}_{\check G}
=\LS^{\reg}_{\check G}\times_{\LS^{\mer}_{\check G}}\LS^{\reg}_{\check G},
\]
它是因子化空间，因为 $\LS^{\reg}_{\check G}$ 与 $\LS^{\mer}_{\check G}$ 都是。于是

\begin{definition}
\label{def:sph-spec}
因子化范畴 $\Sph^{\spec}_{\check G}\in \FactAlgCat$ 为
\[
\IndCoh^{*}\bigl(\LS^{\reg}_{\check G}
\times_{\LS^{\mer}_{\check G}}\LS^{\reg}_{\check G}\bigr).
\]
\end{definition}

然而我们强调，这并不只是把 \cite[附录 A.5]{ABC+24a} 中的 $\IndCoh^{*}$ 构造施于
$\Hck^{\spec}_{\check G}$。事实上，$\IndCoh^{*}$ 在 $\Hck^{\spec}_{\check G}$ 上根本
没有定义。代替它的是等同
\[
\Hck^{\spec}_{\check G}\simeq L^+_\nabla G\backslash L_\nabla G/L^+_\nabla G,
\]
而记号
\[
\IndCoh^{*}\bigl(\LS^{\reg}_{\check G}\times_{\LS^{\mer}_{\check G}}
\LS^{\reg}_{\check G}\bigr)
\]
的定义是
\[
\IndCoh^{*}(L_\nabla G)^{L^+_\nabla G\times L^+_\nabla G}.
\]
（注意 $L_\nabla G$ 是 ind-仿射 ind-概形，故 $\IndCoh^{*}$ 在其上有定义。这与
$\QCoh_{co}$ 在 ind-仿射 ind-概形上有定义的道理相同，见 \cite[附录 A.2]{ABC+24a}。）

它具有由卷积给出的幺半结构，见 \cite[附录 E.6]{ABC+24a}；它按 \cite[附录 E.7]{ABC+24a}
所述作用在 $\Rep(\check G)$ 上。
\srcpage{73}

注意拉回
\[
(\LS^{\mer}_{\check G})_{\Ran_{x_0}}
=\LS^{\mer}_{\check G}\times_{\Ran}\Ran_{x_0}
\]
是因子化 $\LS^{\mer}_{\check G}$-模空间。由于 $\check{\mathfrak{n}}/\check B^{\Ran}_{x_0}$
是因子化 $\LS^{\reg}_{\check G}$-模空间，可知
\[
\check{\mathfrak{n}}/\check B^{\Ran}_{x_0}
\times_{(\LS^{\mer}_{\check G})_{\Ran_{x_0}}}
\check{\mathfrak{n}}/\check B^{\Ran}_{x_0}
\]
是因子化 $\Hck^{\spec}_{\check G}$-模空间。

同样，$\IndCoh^{*}$ 在
$\check{\mathfrak{n}}/\check B^{\Ran}_{x_0}
\times_{(\LS^{\mer}_{\check G})_{\Ran_{x_0}}}
\check{\mathfrak{n}}/\check B^{\Ran}_{x_0}$ 上没有定义，故必须说明记号
\[
\IndCoh^{*}\Bigl(\check{\mathfrak{n}}/\check B^{\Ran_{x_0}}_{x_0}
\times_{(\LS^{\mer}_{\check G})_{\Ran_{x_0}}}
\check{\mathfrak{n}}/\check B^{\Ran_{x_0}}_{x_0}\Bigr)
\]
的含义。然而我们不知道仿射情形有与球面情形中取 $L^+_\nabla G$-双余不变量相平行的定义。
因此把 $\Aff^{\spec}_{\check G}$ 的定义推迟到 3.3.1 小节。

\begin{definition}
\label{def:aff-spec-prelim}
$\FactModCat_{x_0}$ 中的对 $(\Sph^{\spec}_{\check G},\Aff^{\spec}_{\check G})$ 由因子化范畴
$\Sph^{\spec}_{\check G}\in \FactAlgCat$ 与因子化 $\Sph^{\spec}_{\check G}$-模范畴
\[
\Aff^{\spec}_{\check G}
=\IndCoh^{*}\Bigl(\check{\mathfrak{n}}/\check B^{\Ran}_{x_0}
\times_{(\LS^{\mer}_{\check G})_{\Ran_{x_0}}}
\check{\mathfrak{n}}/\check B^{\Ran}_{x_0}\Bigr)
\]
构成。它具有幺半对结构，并按与球面情形类似的方式作用在
$(\Rep(\check G),\QCoh(\check{\mathfrak{n}}/\check B))$ 上。
\end{definition}
